\section{Capítulo~\ref{sec:chemoutput}. La salida química}

Los dos cables hacia el corazón y su distinta velocidad, la realimentación negativa en las cascadas hormonales, el código por frecuencia de porciones, el generador diario y su ajuste por la luz están medidos directamente; el límite $K<8$ es un teorema de la teoría de control deducido en el capítulo; las pulsaciones generadas por la propia realimentación de la cascada son una hipótesis con un experimento sólido a su favor.

{\footnotesize
\setlength{\tabcolsep}{3pt}\renewcommand{\arraystretch}{1.15}
\begin{longtable}{>{\raggedright}p{0.5cm}>{\raggedright}p{4.0cm}>{\raggedright}p{5.6cm}>{\raggedright}p{2.3cm}>{\raggedright\arraybackslash}p{1.7cm}}
\toprule
N.º & Afirmación & Experimento y qué se midió & Fuente & Estado \\
\midrule\endfirsthead
\toprule
N.º & Afirmación & Experimento y qué se midió & Fuente & Estado \\
\midrule\endhead
1 & El marcapasos del corazón late por sí solo unas $100$ veces por minuto; en reposo el freno está pisado y la frecuencia suele ser del orden de $60$--$70$ & Humanos, bloqueo completo de ambos cables hacia el corazón: la frecuencia propia es mayor que la de reposo y disminuye con la edad; en adultos es del orden de $100$ latidos por minuto. Por tanto, en reposo predomina el freno. La frecuencia de reposo de $60$--$70$ es un valor típico, sin cita aparte. & \cite{JoseCollison1970ihr} & medido \\
2 & El acelerador \nt{NORA} y el freno \nt{ACET} son filtros pasa bajos, y el freno es más rápido~\eqref{eq:gasbrake} & Perro, estimulación modulada en frecuencia del nervio vago o simpático cardiaco y función de transferencia hacia la frecuencia auricular: ambas vías son filtros pasa bajos; la simpática tiene una frecuencia de corte mucho menor y un retardo puro añadido de unos $1.7\,\text{s}$. Las características dependen del tono medio, así que la fórmula lineal~\eqref{eq:gasbrake} es una aproximación. & \cite{Berger1989} & medido \\
3 & El freno es rápido porque \nt{ACET} abre el canal de potasio sin segundo mensajero, y \nt{NORA} lo hace mediante una cadena de reacciones & Parches de membrana de células auriculares: el canal de potasio que abre la acetilcolina lo abren las subunidades $\beta\gamma$ de la proteína G unida al receptor, sin segundo mensajero dentro de la célula. La vía de \nt{NORA} a través del AMPc no se cita aquí. & \cite{Logothetis1987gbg} & medido \\
4 & La sangre recorre el cuerpo en un tiempo del orden de un minuto; \nt{ADRE} llega desde el torrente sanguíneo a todos los órganos con receptor & Estimación general de orden de magnitud; así se formula en el capítulo, sin fuente citada. & --- & estimación \\
5 & La cascada hormonal está envuelta en una realimentación negativa: la hormona final inhibe las primeras etapas & Revisión de experimentos: las hormonas de la corteza suprarrenal suprimen la liberación de ACTH por la hipófisis y de la hormona liberadora por el hipotálamo en tres escalas de tiempo: rápida, intermedia y lenta. & \cite{KellerWood1984fb} & medido \\
6 & Tres etapas iguales son estables si $K<8$~\eqref{eq:cascadestable}; en la frontera el periodo es $2\pi\tau/\sqrt3\approx3.6\tau$ & Deducido en el capítulo: a la frecuencia $\sqrt3/\tau$ las tres etapas dan un desfase de $180^\circ$ y una atenuación de $8$ veces. & deducción del capítulo & teoría \\
7 & El error de nivel es $1/(1+K)$, por lo que un regulador así no mantiene el nivel con precisión mejor que un $10\%$ aproximadamente (proposición~\ref{prop:cascade}) & Consecuencia de la fila~6 para realimentación proporcional. El eje real tiene realimentación en varias etapas y en varias escalas de tiempo (fila~5), así que el límite se refiere al modelo~\eqref{eq:cascade} y no está medido. & deducción del capítulo & teoría \\
8 & Con $K=4$ y $7$ el nivel se estabiliza; con $K=12$, pulsaciones regulares con periodo de $3.7$--$3.9\,\tau$ (fig.~\ref{fig:cascade}) & Cálculo con \texttt{models/\allowbreak chem\_\allowbreak models.py}: con $K=12$ los periodos sucesivos son $3.9$, $3.7$, $3.8$, $3.7\,\tau$; con $K=4$ la amplitud al final del cálculo es $0.003$. & \texttt{models/\allowbreak chem\_\allowbreak models.py} & modelo \\
9 & La hormona del estrés se libera en pulsos más o menos una vez por hora; los pulsos los genera la propia realimentación, sin generador aparte & Ratas, infusión constante de hormona liberadora: la ACTH y la corticosterona oscilan con una frecuencia fisiológica, casi horaria; no hace falta una entrada rítmica del hipotálamo para los pulsos. Un modelo hipófisis--suprarrenal con realimentación retardada produce estas oscilaciones. & \cite{Walker2012glc, Walker2010} & hipótesis \\
10 & El destinatario lee la frecuencia de las porciones, no el nivel & Monos sin liberación propia de hormona liberadora de gonadotropinas: la infusión constante no restablece la liberación de hormonas hipofisarias; las porciones una vez por hora sí; volver a la infusión constante suprime de nuevo la respuesta. La fatiga del receptor como filtro pasa altos es una interpretación. & \cite{Belchetz1978} & medido \\
11 & Distintas frecuencias de porciones actúan de forma diferente sobre distintas células de una misma glándula & Los mismos monos: aumentar las porciones a $2$--$5$ por hora reduce ambas hormonas hipofisarias; espaciarlas a una cada $3$~horas reduce una (LH) y aumenta la otra (FSH), de modo que su cociente lo fija la frecuencia. & \cite{Wildt1981gnrh} & medido \\
12 & El ritmo diario lo genera una realimentación negativa intracelular con un retardo de varias horas & Revisión: las proteínas de los genes reloj suprimen con retardo la transcripción de sus propios genes; el lazo de transcripción y traducción da un periodo de alrededor de un día. & \cite{TakahashiJS2017} & medido \\
13 & El generador principal es un grupo de decenas de miles de neuronas que sincroniza el resto del cuerpo & Revisión: el núcleo supraquiasmático es el generador diario principal; sus neuronas aisladas en cultivo son generadores autónomos, y la red las sincroniza; en el ratón, unas $10^4$ neuronas en cada lado. & \cite{Welsh2010scn} & medido \\
14 & El periodo propio en el ser humano es de $24.18$~horas & Personas con luz tenue y un horario desacoplado del día: el periodo de los ritmos de melatonina, temperatura y cortisol es en promedio de $24.18$~horas en jóvenes y mayores, con poca dispersión. & \cite{Czeisler1999} & medido \\
15 & La luz ajusta el generador como un lazo de enganche de fase~\eqref{eq:pll}; hace falta un desplazamiento de unos $11$~min al día & Personas, un solo pulso de luz brillante de $6.7$~horas a distintas horas del día: curva de respuesta de fase de tipo~1 con amplitud de $5$~horas: retraso antes del mínimo de la temperatura corporal, adelanto después. La ecuación~\eqref{eq:pll} es el modelo estándar; $11$~min $=0.18$~horas es aritmética. & \cite{Khalsa2003prc} & medido \\
16 & Sin luz no hay enganche, y el ritmo deriva hacia su periodo propio & Personas totalmente ciegas, sin percepción de luz, que viven en sociedad normal: en cerca de la mitad, los ritmos de melatonina y cortisol siguen su propio periodo, que no coincide con el día. & \cite{Sack1992blind} & medido \\
\bottomrule
\end{longtable}}
